\chapter{更新结论：旧空间不等于当前机会}

\begin{houseviewbox}[核心结论]
截至2026年8月11日盘前、以8月10日收盘为价格基准，历史候选中没有一只同时满足“价格仍处低位、盈利分母可归一、现金可验证、存在当前独立估值锚”四项条件。故本报告\textbf{不发布新的AStock目标价}，也不把券商目标、旧情景价或业绩预告年化为投资结论。
\end{houseviewbox}

\section{决策摘要}

\begin{exhibitbox}[图表1｜8月11日更新后的研究状态]
\small
\begin{tabularx}{\textwidth}{L{1.1cm}L{1.55cm}R{1.05cm}L{2.1cm}X}
\toprule
\textbf{代码} & \textbf{标的} & \textbf{8月10日收盘} & \textbf{更新状态} & \textbf{决定性约束} \\
\midrule
002497 & 雅化集团 & 19.09 & 最高上行监测优先级 & H1仅为预告；FY2025和Q1经营现金流为负；CNY42仅为单一外部锚。 \\
002379 & 宏桥控股 & 19.76 & 事件验证优先级 & H1为预告，CNY29使用调整后利润；等待实际H1、OCF与口径调节。 \\
300014 & 亿纬锂能 & 55.22 & 事件验证优先级 & CNY77与CNY103不可平均；H1收入/OCF未披露且Q1 OCF为负。 \\
000703 & 恒逸石化 & 18.18 & 移出低位清单 & 价格位于一年区间99.96\%，20日上涨37.52\%；实际H1 OCF仍为负。 \\
\bottomrule
\end{tabularx}
\normalsize
\sourcenote{市场数据包、CNINFO公告和归档原始券商PDF；收盘日为2026年8月10日。}
\end{exhibitbox}

\paragraph{排序口径与允许动作。} 图表1的顺序是\textbf{研究投入优先级}，按当前价格位置、下一份实际披露的证据价值、现金/分母缺口和外部锚可审计性综合排列；它\textbf{不是}预期收益排序、投资评级、可投资性排序或仓位建议。雅化、宏桥和亿纬的允许动作仅为事件验证；恒逸的允许动作仅为盈利/现金跟踪。四只原主标的之外，本案还有6只历史高空间样本与3只条件/周期观察样本，均不进入正式目标池。任何证据未达标均按“事件验证 $\rightarrow$ 观察”降级。

\section{与7月结论相比，什么发生了变化}

第一，价格位置重置。恒逸的实际中报利润显著改善，但股价已接近调整后一年高位；因此“盈利兑现”不能再被叙述成“低位重估”。第二，中报信息的质量并不一致：恒逸与江波龙已有实际中报，其余多数只是业绩预告，不能按实际中报或全年利润处理。第三，历史高空间案例没有获得同时满足的当前目标、利润质量与现金证据，故不再以旧模型空间排序。四个原主标的、六个历史高空间样本和三个条件/周期观察样本的边界在第6章统一复核。

\section{读者应如何使用本报告}

本报告是一个研究升级或降级清单：雅化、宏桥、亿纬需要等待下一份实际经营与现金证据；恒逸需要等待新的风险回报重置；其余历史高空间标的只保留监测。它不是交易指令、仓位建议或对外发布的证券评级。
